% BANKS-SOURCE: pdf=255 print=245 kind=appendix id=appendixA
\BanksAppendix{Books}

% BANKS-SOURCE: pdf=255 print=245 kind=prose
This is a brief guide to other books on quantum field theory. The standard modern
textbook is \emph{An Introduction to Quantum Field Theory}, by Peskin and Schroeder \cite{banks-ref-33}.
I recommend especially their wonderful Chapter 5, and all of the calculational sections
between 16.5 and 18.5, as well as Chapters 20 and 21. Every serious student of QFT
should work out the final project on the Coleman--Weinberg potential, which can be
found on page 469. Another standard is Weinberg's three-volume opus \cite{banks-ref-131}.
Here I recommend the marvelous sections on symmetries and anomalies in Volume II. The
technical discussions of perturbative effective field theory are invaluable. The section on
the Batalin--Vilkovisky treatment of general gauge equivalences is also useful. Volume I
should probably be read \emph{after} completing a first course on the subject. It presents an
interesting but idiosyncratic approach to the logical structure of the field. Volume III on
supersymmetry is full of gems. In my opinion, it is flawed by an idiosyncratic notation
and a tendency to obscure relatively simple ideas in an attempt to give absolutely general
discussions. Finally, let me mention a relatively new book by M. Srednicki \cite{banks-ref-168}. I have
not gone through it thoroughly, and I do not agree with the author's ordering of
topics, but the pedagogical style of the sections I have read is wonderful. It is clear that
everyone in the field will turn to this book for all those nasty little details about minus
signs and spinor conventions. I think there is also a chance that it will replace Peskin
and Schroeder as a standard textbook.

I have also enjoyed using the books by Bailin and Love \cite{banks-ref-169} and Ramond \cite{banks-ref-170} in
my many years of teaching the subject. The books by Itzykson and Zuber \cite{banks-ref-171} and
Zinn--Justin \cite{banks-ref-111} are more monographs/encyclopaediae than textbooks, but they con-
tain a wealth of detail on specific subjects in field theory that can be found nowhere
else. I mention especially Zinn--Justin's discussions of the large-order behavior of per-
turbation theory, of the use of field theory for calculating critical exponents, and of
instantons. Other treatments of the field theory/statistical physics interface can be found
in the book by Drouffe and Itzykson \cite{banks-ref-172}, the marvelous book by Parisi \cite{banks-ref-173}, and the
book by Ma \cite{banks-ref-174}. Schwinger's source theory books \cite{banks-ref-175} also belong in the category
of non-textbooks, which contain scads of invaluable information about field theory.

Among older field-theory books, the second volume of Bjorken and Drell \cite{banks-ref-176}
contains lots of useful information, like explicit forms for the spectral representations
for higher spin. The books of Nishijima \cite{banks-ref-177} (Chapters 7 and 8), Bogoliubov and
Shirkov \cite{banks-ref-178}, and even Schweber \cite{banks-ref-179} will reward the really serious student of the
subject.

Finally, I want to mention various shorter documents that I think are essential
reading for students of quantum field theory. The most important is Kogut and

% BANKS-SOURCE: pdf=256 print=246 kind=prose
Wilson \cite{banks-ref-126}, still the best introduction I know of to Wilson's profound ideas about
renormalization. Next is the 1975 Les Houches lecture-note volume \emph{Methods in Field
Theory} \cite{banks-ref-180}, every chapter of which is a gem. Coleman's book \cite{banks-ref-138} is a collection of
lectures on a variety of topics in field theory. I've drawn on it heavily for the material
about instantons and solitons, but the other lectures are also worth reading. The contri-
butions of Adler, Weinberg, Zimmermann, and Zumino to the 1970 Brandeis Summer
School Lectures,\footnote[1]{These lectures also contain a marvelous introduction to string theory by Mandelstam.} and of Weinberg to the 1964 Brandeis volume, are also worth read-
ing \textup{[181]}. Much of Weinberg's material reappears in his textbook \cite{banks-ref-131}. Finally, let me
mention the reprint volume \emph{Selected Papers on Quantum Electrodynamics} \textup{[182]} edited
by Schwinger. The contributions of Feynman and Schwinger in particular should be
read by every student of field theory.

There are lots of other books on field theory, and I apologize to those authors
I haven't mentioned. I've emphasized those texts I've found most useful in my own
career. Others will have different favorites.
